% Chapter 4: Data, models, and experimental setup

\chapter{Data, Models, and Experimental Setup}
\label{ch:setup}

This chapter describes how we produced the answers and the rationales.
Figure~\ref{fig:pipeline} gives the shape of the experiment: the design holds
VCR records and text inputs fixed while varying the image condition and the
presence of polygon overlays or descriptions. Chapter~\ref{ch:framework} defines
the measures applied to the resulting outputs.

\begin{figure}[ht]
\centering
\newcommand{\pipelineheading}{\fontsize{8.85}{10.0}\selectfont}
\newcommand{\pipelinesupport}{\fontsize{6.46}{7.5}\selectfont}
\begin{tikzpicture}[
  font=\pipelineheading,
  >=Latex,
  box/.style={draw=black!55, rounded corners=2pt, align=center,
              inner sep=5pt, minimum height=9mm, fill=black!3},
  wide/.style={box, text width=74mm},
  stage/.style={box, fill=blue!8, draw=blue!45, text width=30mm},
  output/.style={box, fill=orange!10, draw=orange!55, text width=30mm},
  fixed/.style={box, fill=green!8, draw=green!45, text width=24mm}
]
\node[wide] (record) at (1.3,3.3) {\textbf{Shared record}\\
  {\pipelinesupport question, four answer choices, and the fixed description in the description arms}};
\node[wide] (constructor) at (1.3,1.7) {\textbf{Image condition}\\
  {\pipelinesupport source, grey, mismatch, mask, noise, mirror}};

\node[stage] (stage1) at (-2.9,0.0) {\textbf{Stage~1}\\
  {\pipelinesupport own adaptation}};
\node[stage] (stage2) at (3.2,0.0) {\textbf{Stage~2}\\
  {\pipelinesupport own adaptation}};
\node[fixed] (gold) at (7.0,0.0) {gold answer\\
  {\pipelinesupport supplied, unchanged}};

\node[output] (answer) at (-2.9,-1.8) {predicted answer};
\node[output] (rationale) at (3.2,-1.8) {generated rationale};

\draw[->,thick] (record) -- (constructor);
\draw[->,thick] (constructor.south) -- ++(0,-0.35) -| (stage1.north);
\draw[->,thick] (constructor.south) -- ++(0,-0.35) -| (stage2.north);
\draw[->,thick] (gold) -- (stage2);
\draw[->,dashed,black!55] (stage1.east) -- (stage2.west)
  node[midway,fill=white,inner sep=1pt,align=center,font=\pipelinesupport]
  {initialises weights,\\passes no output};
\draw[->,thick] (stage1) -- (answer);
\draw[->,thick] (stage2) -- (rationale);
\end{tikzpicture}
\caption[Two-stage answer and rationale pipeline]{The two-stage evaluation.
The record, question, answer choices, and any description stay the same across
the six image conditions; only the image changes. Stage~1 predicts an answer.
Stage~2 receives the gold answer and generates a rationale for it, so an image
intervention cannot also change the answer Stage~2 must support. Each input arm
has its own Stage~1 and Stage~2 adaptation. Stage~2 starts from the arm's trained
Stage~1 weights, shown by the dashed arrow; no Stage~1 answer output is passed to
Stage~2 in the primary analysis. RQ1 evaluates the answers and RQ2
evaluates the rationales.}
\label{fig:pipeline}
\end{figure}


\section{Dataset and evaluation population}
\label{sec:setup-dataset}

Section~\ref{sec:background-vcr} introduced Visual Commonsense Reasoning (VCR),
its two selection tasks, indexed entities, and region metadata
\citep{DBLP:conf/cvpr/ZellersBFC19}. For this study, each record supplies a movie
frame, a question, four answer choices, the gold answer index, and a human
rationale. The gold answer index is the correctness reference throughout. This
thesis renders an indexed person as a label such as \texttt{person\_0}, and it
uses the polygon form of the region metadata. A given question refers to a
selection of the indexed entities.

The two sets of records come from different VCR splits. Adaptation used 5\% of the
VCR training split: 10,646 records out of 212,923. Evaluation used 10\% of the VCR
validation split: 2,653 records out of 26,534, drawn from 2,400 source frames. The
stored subsets contain unique records, preserve the order of their parent splits,
and are nearly balanced across the four gold answer labels: 2,661, 2,662, 2,662,
and 2,661 for adaptation, and 663, 664, 663, and 663 for evaluation. Both subsets
were constructed with \texttt{data/subset\_balanced\_json.py}. The script grouped a
split's records by their four gold answer labels, shuffled each group with a
generator seeded with 42, and took an equal number from each group, giving the
remainder to one label at a time. It then restored the selected records to
parent-split order. Replaying it reproduces both stored files byte for byte. The
subsets are identified by their filenames and SHA-256
hashes.\footnote{Adaptation used
\texttt{vcr\_train\_5\_pct.json}, SHA-256\\
\texttt{27e156d7043289685ce0723d94fc676610d817b2241097d2f2e70d9b722a0d30}.\\
Evaluation used \texttt{vcr\_val\_10\_pct.json}, SHA-256\\
\texttt{0288740b785021e667c12eb08ab2c222a65f80b2702d598a7ebb7b942b3a5fdc}.} The evaluation set is shared across every arm
and condition. We compared the two files directly: they share no record identifier
and no source frame identifier, so overlap in records and in source
frames between adaptation and evaluation is zero. Several questions can use one frame, so we retain the record count
and the source frame count separately.

Figure~\ref{fig:record-plain-vs-point} shows the two image families on one
record, together with that record's question, choices, and outputs from both
stages. A plain arm uses the unmarked source frame. A point arm uses a copy
specific to the question with magenta polygon outlines and indexed labels on the
referenced regions.

\begin{figure}[ht]
\centering
\begin{tabular}{cc}
\includegraphics[width=0.44\textwidth]{Figures/records/M/source.png} &
\includegraphics[width=0.44\textwidth]{Figures/records/M/polygon.png} \\
\small Source image, used by the plain arms &
\small Polygon overlay, used by the point arms
\end{tabular}

\vspace{5pt}
\begin{minipage}{0.92\textwidth}
\small
\textbf{Question:} What job is person\_0 performing ?\\[2pt]
\textbf{A.} person\_1 is a professional dancer .\\
\textbf{B.} person\_1 is a host .\\
\textbf{C.} person\_0 is a wizard in a play .\\
\textbf{D.} person\_0 is performing as a bartender .\\[2pt]
\textbf{Stage~1 predicted answer:} D\\
\textbf{Gold answer supplied to Stage~2:} D\\
\textbf{Stage~2 final rationale:} person\_0 is standing behind a bar counter .
\end{minipage}
\caption[One VCR record through both stages]{One evaluation record in the two
image forms and through both stages. The plain arms receive the unmarked frame
and the point arms the polygon overlay of the same frame; the question, choices,
and answers are identical in both. Stage~1 predicts an answer, and in the primary
experiment Stage~2 receives the gold answer rather than Stage~1's output. The
displayed prediction and rationale come from the plain arm under the source
condition. The record is illustrative and was not selected for any result.
Original VCR wording and spacing are preserved.}
\label{fig:record-plain-vs-point}
\end{figure}


\section{Four input arms}
\label{sec:setup-arms}

The four arms combine the presence or absence of polygon overlays and fixed
descriptions (Table~\ref{tab:setup-arms}). Each arm was adapted separately, as
Section~\ref{sec:problem-design} explains.

\begin{table}[ht]
\centering
\caption[Four input arms]{Inputs that distinguish the four arms. The VCR
record, question, choices, and answer labels remain matched. Descriptions were
generated from the \mbox{polygon-marked} image in both description arms, so
\texttt{plain\_desc} pairs an unmarked image with text derived from the marked
one.}
\label{tab:setup-arms}
\begin{tabular}{lccL{6.0cm}}
\toprule
Arm & Polygon overlay & Description & Input \\
\midrule
\texttt{plain} & no & no & unmarked source image \\
\texttt{point} & yes & no & image with question-specific polygons \\
\texttt{plain\_desc} & no & yes & source image and fixed description \\
\texttt{point\_desc} & yes & yes & polygon image and the same description \\
\bottomrule
\end{tabular}
\end{table}

We generated descriptions from the polygon-marked images through
\mbox{OpenAI's} API that accepts images, using the \texttt{gpt-4o} model alias. The alias is mutable
and no dated snapshot was pinned, so the exact weights behind it are not
recoverable; Section~\ref{sec:discussion-limitations} records that as a
reproducibility limitation. Appendix~\ref{sec:app-prompts} gives the exact prompt and
generation settings. The description model received the image but not the VCR question, choices,
gold answer, or human rationale as text. The prompt requested a short, factual
scene description using the visible indexed labels. The descriptions are
therefore informed by the annotations, because they were generated from images
containing indexed overlays. This matters most for \texttt{plain\_desc}, which
pairs an unmarked image with a description derived from its marked copy. Where
indexed references occur, the \texttt{plain} arm sees them in the question and
choices but receives no added description linking those labels to the scene.

Each description arm receives the first 256 characters of the stored
description. The cap is applied identically at both stages and in every arm and
condition. We selected 256 characters after trial runs because this setting
avoided GPU out-of-memory errors in the available hardware setup. Because it cuts
at a fixed offset, the supplied text can end in the middle of a word, and
details that appear later in a longer description, including entries for
individual people,
are simply absent from what the model receives. The stored descriptions run to a
median of 445 characters, so 2,624 of the 2,653 records are cut, and 1,674 of them
are cut in the middle of a word. At that
median stored length of 445 characters, the 256-character prefix contains about
58\% of the description.
Appendix~\ref{sec:app-w2c} shows the supplied text exactly as the model
received it for each audited record.

\section{Model and adaptation}
\label{sec:setup-model}

We used Qwen2.5-VL-3B-Instruct
\citep{DBLP:journals/corr/abs-2502-13923}, pinned to model and processor
revision \texttt{66285546d2b821cf421d4f5eb2576359d3770cd3}. We adapted the
model with low-rank adaptation (LoRA)
\citep{DBLP:conf/iclr/HuSWALWWC22}. The base weights remained fixed and were
loaded in four-bit form following QLoRA
\citep{DBLP:conf/nips/DettmersPHZ23}. Computation used bfloat16.

Each arm had one Stage~1 adapter and one Stage~2 adapter. All adaptations used
seed~42, LoRA rank~32, alpha~64, and dropout~0.05, applied to the query, key,
value, and output projections and to the up and down projections of the
feedforward blocks. Stage~2 was trained with the causal language modelling loss on
targets derived from the human rationale. The four Stage~2 checkpoints were trained with
gold answers and fixed for evaluation. Stage~2 begins from the arm's selected
Stage~1 weights and is then trained on the rationale objective; it does not
carry the Stage~1 classifier head, and at inference it receives no Stage~1
hidden state or class score.

Stage~1 predicts through a classifier head. The backbone processes the image
together with the full textual prompt, which contains the question and all four
choices. The Stage~1 representation is
the final hidden state on the text side, taken at the last non-padding token. That vector
passes through a head of the form
\[
  \mathrm{Linear}(H,H)\rightarrow\mathrm{GELU}\rightarrow
  \mathrm{Dropout}(0.1)\rightarrow\mathrm{Linear}(H,4),
\]
where GELU is the Gaussian Error Linear Unit and $H$ is the backbone hidden
size. The four output logits correspond to the
four answer options in their dataset order, and the prediction is the option
with the largest logit. The A--D instruction in the prompt is part of the
textual framing of the task. The classifier gives one explicit score per
candidate answer while leaving the multimodal representation from the backbone
intact.

Table~\ref{tab:setup-training} lists the settings shared by every arm. Both
stages used AdamW with a cosine schedule for the learning rate and a warmup over
the first 5\% of updates. Stage~1 used a learning rate of $1\times10^{-4}$, weight
decay 0.01, gradient accumulation of six steps, and clipping the gradient norm at
1.0. Stage~2 used a learning rate of $5\times10^{-5}$, weight decay 0, and
gradient accumulation of two steps. Both stages used a batch size of two per
device and ran on one device, so accumulation alone determines how many examples
contribute to an update. Stage~1 selected the epoch with the highest validation
accuracy, and for every arm that was epoch~2 of three. Stage~2 ran a fixed budget
of six epochs and retained epoch~6 in every arm. BLEU against the human rationale
was monitored during Stage~2 training, but epochs one to five were scored on a
fixed subsample of 200 records and epoch~6 on all 2,653 records, so those values
from different epochs are not a directly comparable selection curve. The retained
Stage~2 checkpoint was therefore the final checkpoint at the end of its fixed
budget of six epochs. All four arms followed the same training procedure, but
each remained a separately trained system.

The same VCR validation subset of 2,653 records supplied those metrics and the
reported evaluation. The reported results are therefore performance on the
validation set for the selected Stage~1 and retained Stage~2
checkpoints. The study has no separate untouched test set. Every metric taken
at each epoch came from the source condition; no intervention output entered
checkpoint selection or retention. Stage~1 applied an early stopping rule with a
patience of three epochs on validation accuracy; it ran its full budget of three
epochs, so the rule never stopped a run. Stage~2 has no early stopping rule and
ran a fixed budget of six epochs. The runs used
PyTorch~2.9.0, Transformers~4.57.1, PEFT~0.16.0, bitsandbytes~0.48.2, and
Accelerate~1.11.0.

\begin{table}[ht]
\centering
\caption[Shared adaptation settings]{Adaptation settings shared across the four
arms. Retained epochs were read from the run metadata of the eight
checkpoints. Stage~1 kept the epoch with the
highest validation accuracy; Stage~2 kept the final checkpoint of its fixed budget
of six epochs because the BLEU values monitored across epochs did not form a
directly comparable selection curve (Section~\ref{sec:setup-model}).}
\label{tab:setup-training}
\begin{tabular}{lll}
\toprule
Component & Stage~1 & Stage~2 \\
\midrule
Objective & cross entropy over four classes & causal language modelling \\
Role of gold answer & classification target & supplied in prompt \\
Optimiser & AdamW & AdamW \\
Learning rate & $1\times10^{-4}$ & $5\times10^{-5}$ \\
Weight decay & 0.01 & 0 \\
Schedule & cosine; warmup over first 5\% & cosine; warmup over first 5\% \\
Gradient clipping & 1.0 & none \\
Batch size per device & 2 & 2 \\
Gradient accumulation & 6 & 2 \\
Maximum epochs & 3 & 6 \\
Retained epoch, every arm & 2 & 6 \\
Checkpoint retained by & highest validation accuracy & end of fixed budget \\
LoRA rank, alpha, dropout & 32, 64, 0.05 & 32, 64, 0.05 \\
LoRA target modules & \multicolumn{2}{l}{q, k, v, and o projections; feedforward up and down projections} \\
Base weights & \multicolumn{2}{l}{frozen, loaded in four-bit form} \\
Compute precision & \multicolumn{2}{l}{bfloat16} \\
\bottomrule
\end{tabular}
\end{table}

\section{Two-stage prompts}
\label{sec:setup-prompts}

Stage~1 received the image, question, four choices, and optional description.
The text instructed the model to select one of A--D. As
Section~\ref{sec:setup-model} describes, the classifier head's four scores produced
the answer.

The two prompt families label the choices differently. The Stage~1 prompt uses
the letters A--D and the Stage~2 prompt uses zero-based indices (0)--(3), in the
same dataset order. The text and tables a reader sees use the letters
throughout, so index 0 is choice~A, 1 is B, 2 is C, and 3 is D. Indices appear
only in the prompts reproduced in Appendix~\ref{sec:app-prompts} and in
equations.

Stage~2 received the same record and optional description together with a
supplied answer. The prompt conditioned on the gold answer opens by naming the
task, restates the question and the four indexed choices, names the supplied
answer as correct, asks for step-by-step reasoning about the image, and then
gives a literal output template with a \texttt{<reasoning>} and a
\texttt{<final>} field. In the two description arms the first 256 characters of
the stored description are appended after that template, at the very end of the
prompt, under the literal field label \texttt{Image caption:}. The
image is not a token in the prompt text: it is supplied as a separate image item
in the chat template, and the processor inserts its own image tokens.

The two stages label the description field differently. Stage~1 appends the
description under the literal field label \texttt{Caption:}, while primary
Stage~2 uses the literal field label \texttt{Image caption:}. Both carry the
same 256-character description prefix. Appendix~\ref{sec:app-prompts} gives the prompt forms for
Stage~1, Stage~2, and description generation, the ways in which the Stage~2
training and generation prompts differ, and the parser rules that decide which
outputs are eligible.

Stage~2 was trained towards a target in that same two-part format, built from
the VCR human rationale by a formatting heuristic. The rationale is split at
sentence boundaries. When it has more than one sentence, every sentence except
the last becomes the \texttt{<reasoning>} span and the last sentence becomes the
\texttt{<final>} span. When it has exactly one sentence, that sentence is placed
in both spans. The tags therefore reorganise text a human wrote as a single
rationale; they do not add a separately annotated record of hidden reasoning,
and no part of the target was written to describe a model's internal steps. This
is why the primary evaluation unit is the \texttt{<final>} span recovered by the
tolerant parser
of the \emph{generated} text, with the reasoning span and the full generation
reported as sensitivity analyses in
Appendix~\ref{sec:app-rq2-sensitivity}.

Training used the causal language modelling loss over every non-padding token of
the formatted user and assistant sequence. The prompt tokens contributed to the
loss alongside the rationale tokens since the prompt tokens were not masked out.


\section{Six image conditions}
\label{sec:setup-conditions}

The five interventions represent four forms of visual change. Grey removes scene
content from the image. Mask and noise corrupt the source image. Mismatch
replaces it with coherent but unrelated visual content. Mirror changes spatial
orientation while retaining the scene. Source is the reference against which the
other five are compared.

Every record and arm was evaluated under the six conditions in
Table~\ref{tab:setup-conditions}. Questions, choices, supplied gold answers,
and descriptions remained invariant within a record. The mask and noise seeds
depend on dataset index, so construction remains deterministic across shards.

% Two lines shorter, so that the paragraph after Figure~\ref{fig:six-image-conditions}
% begins below the figure instead of being split across it.
\enlargethispage{-2\baselineskip}

\begin{table}[ht]
\centering
\caption[Six image conditions]{Image constructions used at both stages. Each
constructor acts on the image appropriate to the arm.}
\label{tab:setup-conditions}
\begin{tabular}{lL{10.6cm}}
\toprule
Condition & Construction \\
\midrule
source & unmarked image for plain arms or polygon image for point arms \\
grey & a new $256\times256$ RGB image with every channel set to 128; this removes the
scene while still providing an RGB image the processor can accept; it is the only
condition that does not present a frame at its own dimensions \\
mismatch & the target frame of one fixed derangement of the 2,653 records, in the image
family appropriate to the arm; no record maps to itself and no pair shares a source frame \\
mask & one uniform draw in $[0,1)$ per pixel position from NumPy
\texttt{default\_rng}, seeded with the dataset index; every channel at a position is
set to 0 where its draw is below 0.5, so about half of the positions go black \\
noise & independent Gaussian values per colour channel with mean 0 and standard
deviation 80, drawn with a NumPy \texttt{default\_rng} generator seeded with
\texttt{1000000} plus the dataset index, added in 16-bit signed arithmetic,
clipped to 0--255, and cast back to 8-bit RGB \\
mirror & horizontal flip of the image appropriate to the arm \\
\bottomrule
\end{tabular}
\end{table}

Figure~\ref{fig:six-image-conditions} shows these constructors on one
evaluation record.

\begin{figure}[ht]
\centering
\begin{tabular}{ccc}
\includegraphics[width=0.30\textwidth]{Figures/records/C/source.png} &
\includegraphics[width=0.30\textwidth]{Figures/records/C/grey.png} &
\includegraphics[width=0.30\textwidth]{Figures/records/C/mismatch.png} \\
\small Source & \small Grey & \small Mismatch \\[4pt]
\includegraphics[width=0.30\textwidth]{Figures/records/C/mask.png} &
\includegraphics[width=0.30\textwidth]{Figures/records/C/noise.png} &
\includegraphics[width=0.30\textwidth]{Figures/records/C/mirror.png} \\
\small Mask & \small Noise & \small Mirror
\end{tabular}
\caption[Six image conditions for one VCR record]{The six image conditions for
one evaluation record, shown for a plain arm. The question, choices, supplied
answer, and any description stay fixed; only the image changes. Panels share one
canvas so their boxes are comparable and no image is stretched, which is why the
square grey image does not fill its panel. This record is different from the one
in Figure~\ref{fig:record-plain-vs-point} and was selected without consulting
any model output.}
\label{fig:six-image-conditions}
\end{figure}


The mismatch map pairs all 2,653 records one-to-one, contains no fixed points,
and never pairs records that share a source frame. Both stages use the same map,
and it was built without reference to predictions, rationales, or metric
outcomes. The question, answer choices, supplied answer, and description stay as
they are; only the image is replaced, by the target record's image in the form
appropriate to the arm.

The map is reproducible from the evaluation file alone. Load the 2,653 records in
dataset order and seed Python's standard generator with \texttt{20260829}. Copy the
indices 0 to 2,652 and shuffle them, drawing repeatedly from the continuing
generator state, and accept the first permutation in which no index maps to
itself and no source and target share a source frame identifier. The third draw
was the first to satisfy both constraints, and it is the map used at both
stages.

That last detail makes mismatch a slightly different manipulation in the polygon
arms. A \texttt{point} or \texttt{point\_desc} arm receives the target record's
polygon image, so the rendered \texttt{person\_N} labels belong to the target
while the question still refers to the source record's indices. The plain arms
receive an unmarked frame and meet no such conflict. Because the polygon arms
carry labels from the target frame, mismatch comparisons across arms are
descriptive.

Grey differs from the other conditions in a second way that matters for
interpretation. Every evaluation frame is wider than it is tall, between 1,280
and 1,920 pixels across, whereas grey is a fixed 256 by 256 square. Mask, noise,
and mirror keep the source frame's own dimensions, and mismatch presents another
frame at that frame's dimensions. Grey therefore changes the size and shape of
the input as well as removing the scene, so it measures the effect of this
particular replacement rather than of an image without scene content in general.

\section{Inference environments}
\label{sec:setup-lineage}

All Stage~1 predictions, for all six conditions, come from one RTX~3090
environment, with batch size~2 and the same checkpoints and prompts throughout.

Every primary Stage~2 rationale, source and intervention alike, comes from one
environment of two RTX~5090 cards. Generated text can differ across hardware.
Keeping the execution environment fixed means that differences within a source and
intervention pair can be attributed to the image condition rather than the
machine. The primary set contains $4\times6\times2{,}653=63{,}672$
generations.

Stage~2 inference used four-bit NormalFloat (NF4) loading with double
quantisation, bfloat16 compute, greedy decoding, batch size~1, and at most 256
new tokens. Greedy decoding removes sampling. Under the fixed execution setup,
repeated generation reproduced the same text. Identical strings have Drift zero
by definition. The same
model and processor revision, checkpoint, prompt builder, image constructor,
and parser were used across source and intervention conditions.

\section{Gold answer conditioning and secondary regimes}
\label{sec:setup-regimes}

Primary RQ2 uses gold answer conditioning: the gold answer index and text are
identical for a record across all six image conditions. The reason is that
RQ2 asks what the image change does to the rationale. Under \mbox{predicted-answer}
conditioning an image change can also change the Stage~1 answer that Stage~2 is
then asked to justify, so any difference in the text would have two possible
causes at once. Holding the gold answer fixed means that only the image condition
changes within each primary RQ2 comparison.

\subsection{Secondary predicted-answer regime}
\label{sec:setup-pred-a}

PRED-A is a secondary regime that uses only source images. It records what happens
when Stage~2 is given an answer selected by Stage~1 instead of the gold answer,
so it adds no condition to RQ2. It reuses the Stage~2 checkpoints trained with
gold answers and was generated in the RTX~3090 environment. Its results are
reported separately from the primary RQ2 results. All four arms generated with
the same settings: greedy decoding and at most 128 new tokens, half the limit
the primary
generations were given (Section~\ref{sec:setup-lineage}).

The PRED-A prompt differs from the primary generation prompt in four material
ways, and each of them bears on what a comparison could mean. The prompt frames
the task as reflection on the model's own decision. It supplies the Stage~1
answer in place of the gold answer. It adds the instruction \enquote{Continue
until the explanation is complete and well-formed}, which the primary generation
prompt does not carry. In \texttt{plain\_desc} and \texttt{point\_desc} it
appends the description block used during Stage~2 training, which labels the
description as context and instructs the model not to reuse it; the other two
arms receive no description block. A difference between gold and predicted
conditioning therefore cannot be read as the effect of the answer source alone.

The answers came from the Stage~1 validation predictions saved at epoch~3, the
last training epoch. Those predictions reproduce every answer the regime
supplied. The RQ1 results use the retained epoch~2 checkpoint instead, and the
two prediction sets differ on 121 records in \texttt{plain}, 148 in
\texttt{point}, 106 in \texttt{plain\_desc}, and 117 in \texttt{point\_desc},
which is 4.0 to 5.6 per cent of the 2,653 records. PRED-A is therefore not the
configuration that used the retained Stage~1 checkpoint, and its results hold for
the answers actually supplied to those generations.

\subsection{Wrong-answer control}
\label{sec:setup-known-change}

One control changes the supplied answer while keeping the source image fixed. It
provides a known semantic change against which to compare the Drift caused by
image changes. It uses the first 200 records in
deterministic dataset order, giving every arm the same record identifiers in the
same order. Those 200 records come from 185
source frames, and the intervals for the control use the same bootstrap
clustered by source frame as the main analysis. The sample size was fixed before
the final results were generated; it was chosen to keep the additional
generation workload manageable, and no power analysis was made for the number 200.

The control generates one rationale for a predefined alternative answer, giving
800 additional RTX~5090 generations. The alternative
was the Stage~1 source prediction the RQ1 results report, taken from the
retained epoch~2 checkpoint, when that prediction differed from gold; otherwise
it was $(g+2)\bmod4$ for gold index $g$.

\section{Output processing and eligibility}
\label{sec:setup-output-processing}

Stage~2 requested a \texttt{<reasoning>} span and a shorter \texttt{<final>}
span. We stored the raw generation and both strict and tolerant parser outcomes.
The tolerant parser accepts differences in case and spacing and tag names beginning
with \texttt{reason} or \texttt{final}, including the malformed variants documented
in Appendix~\ref{sec:app-prompts}. When no usable tag is present it recovers the
last sentence under one fixed fallback rule. That appendix states the rules exactly.

The primary rationale unit is a final rationale span containing text, recovered by
the tolerant parser. A comparison relative to source is eligible only when both the
source and condition spans are usable for the same record. Missing text was not
replaced with the reasoning span, raw output, text from another condition, or
content generated by the model. Chapter~\ref{ch:results} reports raw, usable, and
pairable counts before the Drift results, and
Figure~\ref{fig:rq2-paired-example} shows one record's spans under three image
conditions.
